\chapter*{ABSTRACT}
\addcontentsline{toc}{chapter}{ABSTRACT}
\doublespacing % Anna University Rule: Double line spacing for Abstract
\normalsize
{\Large \textbf{ABSTRACT} \par}
\vspace{1cm}
Traditional staff attendance and monitoring systems often rely on manual processes, facial recognition, or GPS-based tracking. However, these approaches face significant limitations, particularly in indoor environments where GPS signals are unreliable and facial recognition accuracy is affected by lighting conditions, camera quality, and user positioning. To address these challenges, this project proposes a context-aware indoor staff presence verification system using Bluetooth Low Energy (BLE). In this system, each staff member is assigned a BLE-enabled identification card that continuously broadcasts a unique identifier. BLE receivers, implemented using ESP32 microcontrollers and strategically placed within classrooms, detect these signals and measure the Received Signal Strength Indicator (RSSI) to estimate proximity. 

The system verifies staff presence by correlating detected location data with predefined classroom assignments and scheduled timetables, ensuring that attendance is recorded only when the staff member is present in the correct location at the appropriate time. The proposed system integrates a centralized database and a web-based interface developed using modern technologies such as React.js, Express.js, and MongoDB. It includes modules for staff registration, BLE-based detection, and attendance management. Once verified, attendance is recorded automatically, capturing entry and exit times without requiring manual intervention. Additionally, the system generates real-time alerts for late arrivals, absences, or incorrect classroom presence, enabling efficient monitoring and administrative control. Compared to existing systems, this solution offers improved accuracy, reliability, and efficiency in indoor environments. It is cost-effective, energy-efficient, and scalable, making it suitable for deployment in smart campus infrastructures.
